An automorphism acts as $A \otimes B$ with $(A, B) \in U(m) \times U(k)$, and the kernel of the action $(A, B) \mapsto A \otimes B$ is the set of pairs with $A \otimes B = 1$. Such a pair has both factors scalar: $Av \otimes Bw = v \otimes w$ for all $v$, $w$ forces $Av$ proportional to $v$ for every $v$, so $A = \lambda I_{m}$ (a linear map fixing every line is a scalar), and then $B = \lambda^{-1}I_{k}$; so the kernel is $\{(\lambda I_{m}, \lambda^{-1}I_{k})\}$, a single $U(1)$. On Lie algebras, $\uu(m) \oplus \uu(k) = \su(m) \oplus \su(k) \oplus \uu(1) \oplus \uu(1)$, and quotienting by the diagonal $\uu(1)$ kernel leaves $\su(m) \oplus \su(k) \oplus \uu(1)$. For $m = k = 3$ this is $\su(3) \oplus \su(3) \oplus \uu(1)$, dimension $8 + 8 + 1 = 17$, not the Standard Model's twelve. For $m = 2$, $k = 3$ it is \eqref{eq:algebra}, dimension $3 + 8 + 1 = 12 = 4 + 9 - 1$. The script confirms both dimensions as exact ranks: the action of an explicit thirteen-element basis of $\uu(2) \oplus \uu(3)$ on $\C^{2} \otimes \C^{3}$ has image of rank twelve, and that of the eighteen-element basis of $\uu(3) \oplus \uu(3)$ on $\C^{3} \otimes \C^{3}$ has image of rank seventeen.
